% Revised blocks for the thesis (English).
% Replace the corresponding sections in the main Overleaf file.

% ===== ABSTRACT =====
\begin{abstract}
\addchaptertocentry{\abstractname}
Population ageing requires territorial responses in housing, mobility, health,
social participation, digital inclusion, and support for caregivers. Relevant
initiatives are documented by public institutions, associations, and local
authorities, but the resources are scattered across heterogeneous web pages,
PDF documents, and action sheets. This thesis presents \textit{Index\_Insight},
a reproducible pipeline that transforms this dispersed documentation into a
structured corpus for exploratory analysis.

The corpus contains 1,153 documents from six sources: CNAV, Fondation de
France, HCFEA, IReSP Autonomie, the SIAGE action-sheet corpus, and the
Francophone Age-Friendly Cities and Communities network (VADA). The pipeline
combines source-specific collection, text normalisation, relevance scoring,
structured extraction, and quality-control outputs. An initial rule-based
pipeline was progressively complemented by local LLM-assisted extraction using
Qwen2.5 through Ollama, with constrained JSON generation and controlled
metadata values.

The work delivers a documented corpus and analytical outputs describing themes,
initiative types, stakeholders, territories, and temporal information. The
evaluation distinguishes metadata completion from factual correctness: the
former is measured automatically, whereas the latter requires human reference
annotations. The thesis discusses the benefits and limitations of local
LLM-assisted extraction for heterogeneous territorial data, and identifies the
conditions required before the corpus can support a public search interface or
decision-support use.
\end{abstract}

% Add at the end of "Objectives of the thesis".
\paragraph{Research question.}
This internship addresses the following research question: \emph{to what
extent can a locally deployed language model transform heterogeneous textual
descriptions of territorial initiatives related to population ageing into
reliable, structured, and reusable metadata?} Three operational questions
follow: (i) how can documents with different formats and editorial structures be
collected and normalised in a traceable way; (ii) which metadata can be extracted
under controlled vocabularies; and (iii) how should corpus quality be assessed
without confusing field completion with factual correctness?

% ===== REPLACE THE CORPUS/METHODOLOGY SECTIONS =====
\section{Data collection and corpus construction}

The corpus was constructed from documents describing concrete initiatives,
programmes, research projects, or action sheets related to ageing, autonomy,
social inclusion, housing, mobility, health, participation, or support for
caregivers. The unit of analysis is one documentary record: one web page, PDF
description, or action sheet stored as one text file. The collection process
combined an internal corpus supplied by the SIAGE Chair and publicly accessible
institutional or associative resources.

For each external source, a dedicated collection script was used because page
structures, navigation mechanisms, and available metadata differed. The scripts
retrieve textual content and, where appropriate, downloadable PDF material.
Only documents containing a substantial description of an initiative or project
were retained. Navigation pages, duplicated pages, empty pages, and pages
without a meaningful connection to territorial adaptation to ageing were
excluded whenever they could be identified automatically or during inspection.

Table~\ref{tab:final-corpus} reports the corpus actually processed in the
project. It supersedes exploratory source lists: CNSA and Fondation Médéric
Alzheimer were considered during source exploration but are not represented as
separate final source directories in the delivered corpus.

\begin{table}[H]
\centering
\begin{tabular}{lrp{7cm}}
\toprule
Source & Documents & Nature of the material \\
\midrule
CNAV & 42 & Institutional documents and prevention-related initiatives. \\
Fondation de France & 43 & Associative or territorial initiative descriptions. \\
Fiches actions SIAGE & 60 & Internal action sheets supplied by the host institution. \\
HCFEA & 23 & Institutional reports and ageing-related public-action resources. \\
IReSP Autonomie & 131 & Research-project descriptions in the field of autonomy. \\
VADA & 854 & Descriptions of local age-friendly initiatives. \\
\midrule
Total & 1,153 & Heterogeneous textual corpus. \\
\bottomrule
\end{tabular}
\caption{Final corpus used in the project.}
\label{tab:final-corpus}
\end{table}

The corpus is not a statistically representative sample of all French ageing
initiatives. VADA contributes 854 of the 1,153 records, or approximately
74\%. This imbalance must be considered when interpreting aggregate statistics:
frequent themes or initiative types may partly reflect the coverage of this
source rather than the full French landscape. Raw, cleaned, intermediate, and
structured files are retained separately to preserve provenance and enable
inspection or reprocessing.

\section{Cleaning, relevance scoring, and document preparation}

The collected texts contain formatting artefacts arising from HTML extraction,
PDF conversion, line breaks, duplicated headers, and inconsistent spacing.
Cleaning operations normalise whitespace and line endings, remove obvious
presentation noise, standardise file names, and preserve one document per file.
The objective is not to erase linguistic content, but to create stable textual
inputs for inspection and extraction.

An exploratory relevance score based on weighted ageing-related keywords was
also implemented. It gives priority to specific notions such as \textit{VADA},
\textit{ageism}, or \textit{age-friendly city}, and lower weights to broader
terms such as \textit{autonomy} or \textit{housing}. Scores are normalised by
document length and help prioritise the inspection of unusual documents. This
score is a quality-control tool, not a supervised classifier or a statistically
validated measure of relevance.

\section{Processing architecture and structured extraction}

The project uses a modular Python architecture. Each major transformation has
its own script and output directory, preserving intermediate artefacts and
facilitating debugging. The initial pipeline cleaned documents, segmented text,
detected sections, extracted candidate values, and assigned heuristic confidence
scores. This approach was transparent but difficult to generalise across
documents with very different structures. It was therefore retained as an
experimental and quality-control baseline rather than treated as the sole final
solution.

The project then introduced constrained LLM-assisted extraction. The current
implementation calls the local model \texttt{qwen2.5:1.5b} through the Ollama
chat API. Local inference was selected to avoid external API dependence, retain
control over the workflow, and keep iterative experimentation affordable. This
model identifier should be harmonised throughout the manuscript if another model
was ultimately used to generate the delivered database.

For each document, the script supplies a structured prompt and requests only
valid JSON. The current VADA extraction script uses temperature 0, a context
window of 8,192 tokens, a maximum prediction length of 700 tokens, and an input
truncation threshold of 2,500 characters. Truncation is a practical safeguard,
but it can omit relevant information located late in a document. Raw model
responses are retained to support later inspection.

The expected JSON schema contains \texttt{titre}, \texttt{territoire},
\texttt{echelle}, \texttt{public\_vise}, \texttt{nature\_initiative},
\texttt{description\_generale}, \texttt{contexte}, \texttt{problematique},
\texttt{solution\_envisagee}, \texttt{objectifs},
\texttt{porteur\_initiative}, \texttt{date}, \texttt{thematique}, and
\texttt{source\_site}. Empty strings or lists are required when a value cannot
be supported by the document. Controlled vocabularies are imposed for target
audience and territorial scale; concise, homogeneous labels are requested for
themes and initiative types.

\begin{center}
Raw source $\rightarrow$ text extraction $\rightarrow$ cleaning and inspection
$\rightarrow$ constrained prompt $\rightarrow$ local LLM $\rightarrow$
validated JSON $\rightarrow$ descriptive statistics and manual review.
\end{center}

% ===== REPLACE THE IMPLEMENTATION CHAPTER =====
\chapter{Implementation}

\section{Development environment and reproducibility}

The pipeline was developed in Python in an isolated virtual environment. Its
dependencies cover HTTP collection, HTML parsing, PDF extraction, JSON and CSV
processing, spreadsheet production, and plotting. LLM inference is provided
locally by Ollama. Reproducibility requires recording the Python version,
package versions, Ollama version, model identifier, execution date, prompt
version, and inference parameters.

The core workflow does not rely on an external commercial LLM API. However,
reproducibility is not perfect across machines: a model update, changed source
page, or modified prompt can alter outputs. Raw input text, generated JSON, and
archived raw LLM answers are therefore essential provenance artefacts.

\section{Main modules and outputs}

Data-collection scripts retrieve relevant pages, extract titles and content,
download or parse PDF material when necessary, and write one text file per
resource. Cleaning and file-name utilities prepare the raw material. The
rule-based pipeline contains stages for cleaning, segmentation, section
detection, candidate extraction, scoring, JSON consolidation, and review-table
generation. The LLM module builds a French prompt, submits it to Ollama, checks
or repairs superficial JSON-formatting issues, and saves both structured and raw
responses.

The principal output is a collection of JSON records usable by downstream
software. Complementary outputs include a consolidated JSON file, a CSV review
table, Excel workbooks, and charts describing themes, initiative types,
stakeholders, territories, temporal mentions, and comparisons between SIAGE
action sheets and other documents. They are exploratory outputs, not estimates
of the real-world frequency of initiatives in France.

\section{Difficulties and implemented solutions}

Structural heterogeneity was the main difficulty: a heading in one source can
be a paragraph or be absent in another. The modular pipeline and separate raw
and intermediate directories made this variability inspectable. Constrained LLM
extraction reduced the need for a large number of source-specific rules, while
controlled vocabularies limited output variation. Duplicated boilerplate,
incomplete dates, ambiguous territorial references, and text-encoding artefacts
remain issues requiring cleaning, harmonisation, and manual review.

% ===== REPLACE THE EVALUATION CHAPTER =====
\chapter{Evaluation}

\section{Evaluation objectives}

The evaluation has two distinct objectives. The first measures technical
usability: records should be valid, fields should be populated, and provenance
should be retained. The second measures factual extraction quality: a field must
be correct and supported by the source text. These objectives must not be
confused. A complete record can still contain an incorrect or overly inferred
value.

\section{Automatic coverage results}

The rule-based consolidated database contains 1,153 records. Its review table
measures field completion and an internal heuristic confidence score. The
results in Table~\ref{tab:coverage-results} demonstrate high field population,
but neither measure is a human-judged accuracy score nor a measure of LLM
performance.

\begin{table}[H]
\centering
\begin{tabular}{lrrr}
\toprule
Source & Records & Mean completion rate & Mean heuristic confidence \\
\midrule
CNAV & 42 & 0.99 & 0.58 \\
Fiches actions SIAGE & 60 & 0.94 & 0.53 \\
Fondation de France & 43 & 0.97 & 0.58 \\
HCFEA & 23 & 0.97 & 0.54 \\
IReSP Autonomie & 131 & 0.98 & 0.63 \\
VADA & 854 & 1.00 & 0.60 \\
\midrule
Overall & 1,153 & 0.99 & 0.59 \\
\bottomrule
\end{tabular}
\caption{Automatic coverage of the rule-based consolidated database.}
\label{tab:coverage-results}
\end{table}

Completion covers title, territory, scale, target audience, stakeholder, date,
theme, source, and five descriptive fields. The confidence figure derives from
rule-based candidate scores. It is useful for prioritising manual review, not as
a calibrated probability that the field is correct.

\section{Human reference evaluation protocol}

To evaluate LLM-assisted extraction rigorously, a manually annotated reference
set should be created before final claims about accuracy are made. The sample
should be stratified by source, including smaller sources, and annotators should
receive a guideline defining each field, acceptable inference, and missing
information. For categorical fields (\texttt{public\_vise}, \texttt{echelle},
\texttt{nature\_initiative}, and \texttt{thematique}), report precision,
recall, and macro F1. For free-text fields, ask human judges whether the output
is supported by the text, useful, and free from unsupported information. Record
JSON validity, empty-value correctness, and time per document. A second
annotator on a subset would allow reporting agreement or Cohen's $\kappa$.

\section{Results analysis and limitations}

The automatic results show that the pipeline creates nearly complete structured
records across all six sources, but the mean heuristic confidence remains
moderate (0.53 to 0.63 by source). This confirms that coverage alone is
insufficient: territory, stakeholder, context, problem statement, and solution
may be populated while remaining ambiguous or weakly supported. The charts and
spreadsheets are valuable corpus-exploration tools, but the predominance of VADA
documents prevents interpreting them as estimates for the whole French context.

The main limitation is the absence of a fully annotated gold standard. LLM
outputs can also vary with prompt wording, model version, truncation, and source
quality. The current results establish feasibility and coverage; human reference
evaluation is still necessary to establish field-level accuracy.

% ===== REPLACE DISCUSSION, CONCLUSION, AND APPENDIX CONTENT =====
\chapter{Discussion}

\section{Contributions of the internship}

This internship delivers a documented corpus of 1,153 resources on territorial
adaptation to ageing; a traceable architecture preserving raw, intermediate,
and structured data; a constrained local LLM extraction workflow complementing
a transparent rule-based baseline; and analytical outputs for exploratory
comparison between initiatives, themes, stakeholders, and territorial scales.

\section{Ethical, legal, and methodological considerations}

The corpus consists of publicly accessible institutional and associative
resources and is intended for documentary analysis. Collection must nevertheless
respect source conditions of access, website terms, and requests not to overload
servers. Provenance should be retained whenever a document is reused. The system
should not make consequential decisions about individuals or territories. Human
validation is necessary before model-extracted information is published as fact
or used in policy analysis.

\section{Limitations and future work}

The corpus is unbalanced across sources and inherits their editorial choices.
Text quality varies, and web or PDF extraction can preserve boilerplate,
incomplete metadata, and encoding artefacts. Territorial references may describe
a location, organisation, or project area; initiative holders may be confused
with partners. Controlled vocabularies reduce variation but can hide meaningful
distinctions. Future work should prioritise a human-annotated benchmark,
taxonomy harmonisation, duplicate detection, explicit provenance fields,
comparison of prompts or local models, and then faceted or hybrid semantic
search with an interface designed with INSIGHT users.

\chapter{Conclusion}

This thesis presented \textit{Index\_Insight}, a pipeline for structuring
heterogeneous French-language documentation on territorial responses to
population ageing. Starting from 1,153 documents from six sources, the project
combines collection, cleaning, exploratory relevance scoring, structured
extraction, provenance preservation, and descriptive analysis. The transition
from rule-based extraction to constrained local LLM extraction was motivated by
the diversity of sources and the need for a scalable workflow under local
control. The project demonstrates that a rich structured corpus can be created
from dispersed documentary resources, while also showing that automatic
completion must not be confused with factual accuracy. A human reference
evaluation is the next essential step before the corpus supports robust research
or public-facing retrieval tools.

\appendix
\chapter{Appendices}
\section{Recommended appendix contents}
\begin{itemize}
\item Full version and date of each data-collection script.
\item Final LLM prompt, model identifier, and inference parameters.
\item JSON schema and controlled vocabularies.
\item One raw source text, its generated JSON, and human review.
\item Annotation guideline and sampling procedure for the reference set.
\item Additional charts and source-level descriptive tables.
\end{itemize}

